\documentclass[11pt]{article}
\usepackage[a4paper,margin=27mm]{geometry}
\usepackage{amsmath,amssymb,mathtools}
\usepackage{fontspec,unicode-math}
\setmainfont{Latin Modern Roman}
\setmathfont{Latin Modern Math}
\usepackage{xurl,hyperref}
\hypersetup{colorlinks=true,urlcolor=blue,linkcolor=blue}
\newcommand{\sref}[2]{\hyperlink{source:#1:#2}{[S#2]}}
\newcommand{\cataloguescope}[1]{\paragraph{Mathematical question.} #1}
\begin{document}
\section{O22a — Discrete logarithms modulo a composite}\label{statement}

\subsection{Definitions and mathematical statement}

\paragraph{Shared notation.}$\mathbb N=\{1,2,\ldots\}$, and $\mathbb Z,\mathbb Q,\mathbb R,\mathbb C$ denote the integers, rationals, reals and complex numbers. Any other conventions are specified locally.


Integers are supplied in binary; polynomial time means a polynomial number of classical bit operations in total input length. Given positive integers $g,b,n$ with $n\ge2$, seek a positive integer exponent $x$ satisfying $g^x\equiv b\pmod n$. The computational question is a search problem promised a solution; no coprimality assumption on $g$ is imposed in C22.
\paragraph{Statement recorded in the supplied source.}
Is there a deterministic polynomial-time algorithm which produces such an $x$ on every input for which a solution exists? \sref{statement}{2}
\subsection{Short English statement}
Can a modular discrete logarithm be found deterministically in polynomial time for an arbitrary modulus, including composites, whenever it exists?
\subsection{Sources}
\begin{description}
\item[\hypertarget{source:statement:1}{S1}] ULAM.AI and the UnsolvedMath Contributors, UnsolvedMath: A Curated Collection of Open Mathematics Problems, catalogue entry. CC BY 4.0. \url{https://huggingface.co/datasets/ulamai/UnsolvedMath}.
\item[\hypertarget{source:statement:2}{S2}] Leonard M. Adleman and Kevin S. McCurley, Open Problems in Number Theory Complexity, II, in Algorithmic Number Theory, LNCS 877 (1994), pp. 291--322, C22/O22a, manuscript p. 17 (LNCS p. 307), and Definition 2. \url{https://mccurley.org/papers/open.ps.gz}.
\end{description}
\label{end:statement}
\end{document}
